%% Sommes partielles S_N, points discrets (pas de trait continu, pour
%% bien montrer que S_N est une suite) : sum 1/k (diverge -- affiche
%% jusqu'a N=60 pour bien voir que ca ne plafonne pas) vs sum 1/k^2
%% (converge vers pi^2/6, ligne pointillee). Consomme par slides.
\begin{tikzpicture}[x=0.15cm,y=0.9cm]
    \draw[-Latex] (-1,0) -- (62,0) node[right] {\small $N$};
    \draw[-Latex] (0,-0.3) -- (0,5.4) node[above] {\small $S_N$};
    \node at (30,5.8) {\small $\sum 1/k$ diverge};
    \pgfmathsetmacro{\Ssum}{0}
    \foreach \k in {1,...,60}{
        \pgfmathparse{\Ssum+1/\k}
        \xdef\Ssum{\pgfmathresult}
        \fill[niceblue] (\k,\Ssum) circle (1.3pt);
    }
\end{tikzpicture}
\hspace{1.6cm}
\begin{tikzpicture}[x=0.15cm,y=1.7cm]
    \draw[-Latex] (-1,0) -- (62,0) node[right] {\small $N$};
    \draw[-Latex] (0,-0.15) -- (0,2.05) node[above] {\small $S_N$};
    \node at (30,2.2) {\small $\sum 1/k^2 \to \pi^2/6$};
    \draw[dashed,gray] (0,{pi*pi/6}) -- (60,{pi*pi/6});
    \pgfmathsetmacro{\Ssum}{0}
    \foreach \k in {1,...,60}{
        \pgfmathparse{\Ssum+1/(\k*\k)}
        \xdef\Ssum{\pgfmathresult}
        \fill[red] (\k,\Ssum) circle (1.3pt);
    }
\end{tikzpicture}
